% answers_short.tex —— 2014 年·湖南卷·理科 · 简短答案速查
% 由 tools/make_short_answers.py 自 answers.tex 的 \gksolhead/\gkanswer 提取，
% 并与 2014/2014_数学真题_湖南理_答案要点.md 逐题比对后生成，勿手改（改源 MD 后重跑）。

\begin{gknotice}
\textbf{使用说明}：本节只给最终答案便于快速核对，完整推导见解析版。答案为编写者独立推导，\textbf{非官方标准答案}。
\end{gknotice}

\gkpart{选择题}
\begin{center}\begin{tabular}{ccccccccccc}
\toprule
题号 & 1 & 2 & 3 & 4 & 5 & 6 & 7 & 8 & 9 & 10 \\
\midrule
答案 & B & D & C & A & C & D & B & D & A & B \\
\bottomrule
\end{tabular}\end{center}

\gkpart{填空题}
\begin{itemize}[leftmargin=2.6em]
  \item[\textbf{11.}] $\rho(\cos\theta-\sin\theta)=1$（即 $\rho\cos\left(\theta+\dfrac\pi4\right)=\dfrac{\sqrt2}2$）。曲线是圆心 $(2,1)$、半径 $1$ 的圆，$|AB|=2$ 说明 $l$ 过圆心，故 $l:y=x-1$，化为极坐标即得。
  \item[\textbf{12.}] $\dfrac32$。设 $M$ 为 $BC$ 中点，则 $A$，$M$，$O$ 共线，$BM=\sqrt2$，$AM=\sqrt{3-2}=1$；由 $R^2=(R-1)^2+2$ 得 $R=\dfrac32$。
  \item[\textbf{13.}] $-3$。$|ax-2|<3$ 即 $-1<ax<5$，与 $\left(-\dfrac53,\dfrac13\right)$ 比较端点并由 $a<0$ 得 $a=-3$。
  \item[\textbf{14.}] $-2$。可行域顶点为 $(k,k)$，$(4-k,k)$，$(2,2)$，目标值依次为 $3k$，$8-k$，$6$；最小值为 $-6$ 只有 $3k=-6$，故 $k=-2$。
  \item[\textbf{15.}] $1+\sqrt2$。取 $O$ 为原点、$AD$ 为 $x$ 轴，得 $C\left(\dfrac a2,-a\right)$，$F\left(\dfrac a2+b,b\right)$，由 $C$ 在抛物线上得 $p=a$，再代入 $F$ 得 $b^2=a^2+2ab$，故 $\dfrac ba=1+\sqrt2$。
  \item[\textbf{16.}] $\sqrt7+1$。$\overrightarrow{OA}+\overrightarrow{OB}=(-1,\sqrt3)$，所求向量为 $(2+\cos t,\sqrt3+\sin t)$，其模不超过 $\sqrt{2^2+3}+1=\sqrt7+1$，且可取等号。
\end{itemize}

\gkpart{解答题}
\begin{itemize}[leftmargin=2.6em]
  \item[\textbf{17.}] (1) $\dfrac{13}{15}$。用对立事件：$1-\left(1-\dfrac23\right)\left(1-\dfrac35\right)=1-\dfrac2{15}=\dfrac{13}{15}$；(2) $E(X)=140$（万元）。$X$ 取 $0$，$100$，$120$，$220$，概率依次为 $\dfrac2{15}$，$\dfrac15$，$\dfrac4{15}$，$\dfrac25$，期望为 $100\cdot\dfrac15+120\cdot\dfrac4{15}+220\cdot\dfrac25=140$
  \item[\textbf{18.}] (1) $\dfrac{2\sqrt7}7$。由余弦定理 $\cos\angle CAD=\dfrac{1+7-4}{2\sqrt7}=\dfrac{2\sqrt7}7$；(2) $3$。$\sin\angle BAC=\sin\left(\angle BAD-\angle CAD\right)=\dfrac{\sqrt3}2$，即 $\angle BAC=60^\circ$，再由正弦定理 $BC=\dfrac{\sqrt7\sin60^\circ}{\sin\angle CBA}=3$
  \item[\textbf{19.}] (1) 见解析版；(2) $\dfrac{2\sqrt{57}}{19}$。建系求两平面的法向量，二面角余弦为 $\dfrac{2\sqrt{57}}{19}$
  \item[\textbf{20.}] (1) $p=\dfrac13$。递增给出 $a_{n+1}-a_n=p^n$，$a_2=1+p$，$a_3=1+p+p^2$，由 $4a_2=a_1+3a_3$ 得 $3p^2-p=0$，舍 $p=0$ 得 $p=\dfrac13$；(2) $a_n=\dfrac43+\dfrac{(-1)^n}{3\cdot2^{n-1}}$（$n\in\mathbf N^*$）。两子列单调迫使 $a_{n+1}-a_n=\dfrac{(-1)^{n+1}}{2^n}$，求和即得
  \item[\textbf{21.}] (1) $C_1:\dfrac{x^2}2+y^2=1$；$C_2:\dfrac{x^2}2-y^2=1$。由 $e_1e_2=\dfrac{\sqrt3}2$ 得 $b^2=\dfrac{a^2}2$，再由 $|F_2F_4|=\dfrac a{\sqrt2}(\sqrt3-1)=\sqrt3-1$ 得 $a=\sqrt2$；(2) $2$。设 $AB:x=my-1$，得 $|AB|=\dfrac{2\sqrt2(1+m^2)}{m^2+2}$，$|PQ|=\dfrac{2\sqrt{m^2+4}}{\sqrt{2-m^2}}$，面积 $S=\dfrac{2\sqrt2\sqrt{1+m^2}}{\sqrt{2-m^2}}$，在 $m=0$ 处取最小值 $2$
  \item[\textbf{22.}] (1) $a\geqslant1$ 时在 $(0,+\infty)$ 上单调递增；$0<a<1$ 时在 $\left(0,2\sqrt{\dfrac{1-a}a}\right)$ 上单调递减、在 $\left(2\sqrt{\dfrac{1-a}a},+\infty\right)$ 上单调递增。因 $f'(x)=\dfrac{ax^2+4a-4}{(1+ax)(x+2)^2}$；(2) $\left(\dfrac12,1\right)$。两个极值点要求 $0<a<1$ 且 $a\neq\dfrac12$，此时 $f(x_1)+f(x_2)=\ln\left[(2a-1)^2\right]+\dfrac{4(1-a)}{2a-1}>0$ 等价于 $a>\dfrac12$
\end{itemize}
